\ifx\danslelivre\undefined
\RequirePackage{luatex85}
\documentclass[10pt]{../fichiers-configuration-latex/smfart}
\usepackage{polyglossia}
%\usepackage[francais]{babel}
\usepackage{fontspec}
\usepackage[math-style=ISO,bold-style=upright,sans-style=upright]{unicode-math}
%\usepackage[mathletters]{ucs}
%\usepackage[utf8]{inputenc}
%\usepackage[T1]{fontenc}
\usepackage[unicode,colorlinks=true,linkcolor=blue,citecolor=Sepia,urlcolor=magenta]{hyperref}
\setmainlanguage{french}
\setmainfont[Ligatures=TeX]{Linux Libertine O}
% variante possible : TeXGyrePagella ?
\setmathfont{XITS Math}

% maths
\usepackage{amsmath}
%\usepackage{amssymb}
%\usepackage{mathrsfs}
%\usepackage{stmaryrd}
%\usepackage{euscript}

% couleurs
\usepackage[usenames,dvipsnames]{xcolor}

% graphiques
\usepackage{graphicx}
\usepackage[all]{xy}
% Xy-pic 3.8.9's PDF driver assumes that \fontdimen8\textfont3 is a
% legacy TeX rule width.  Under LuaTeX with unicode-math/XITS Math it is
% instead 6.62pt, which turns every arrow shaft into a thick black bar.
% Override only xypdf's private cached width; do not alter the math font.
\makeatletter
\AtBeginDocument{%
  \@ifundefined{xP@preclw}{}{%
    \edef\xP@preclw{0.4pt}%
    \edef\xP@lw{\xP@dim\xP@preclw}%
    \PackageInfo{sga4-xypdf-compat}{Xy-pic PDF line width set to \xP@preclw}%
  }%
}
\makeatother
%\usepackage{tikz}
%\usetikzlibrary{matrix,arrows,calc}
%\tikzset{column sep=3em,row sep=2em,text height=1.5ex, text depth=0.25ex,>=angle 90}
%\tikzset{description/.style={fill=white,inner sep=1pt}}

% hyperliens, 'etiquettes
\usepackage{xr-hyper}
\usepackage{../fichiers-configuration-latex/smfthm}
%\ifx\cacheretiquettes\undefined
%\usepackage[inline]{showlabels} \else \fi % pour rendre visibles les 'etiquettes
%\usepackage{../fichiers-configuration-latex/makeidx} % pour faire un index
%\usepackage[pagebackref,unicode]{hyperref}  % pour voir o`u sont cit'ees les r'ef'erences
\hypersetup{unicode,colorlinks=true,linkcolor=blue,citecolor=Sepia,urlcolor=magenta,
  hypertexnames=false}
%\hypersetup{breaklinks=true} %permet le retour `a la ligne dans les liens trop longs. Incompatible avec pdflatex
%\renewcommand*{\backref}[1]{{\small [cit\'e page(s)~#1.]}}

\usepackage{xspace}
\usepackage{pifont} % pour expose XVII

\renewcommand{\thefootnote}{\roman{footnote}} % pour mettre les notes en bas de page sous la forme (ii) et non (2), qui ressemble trop (?) a un exposant.

\newcounter{pageoriginal}
\newcommand{\pageoriginale}{\refstepcounter{pageoriginal}\marginpar{\footnotesize\textbf{\thepageoriginal}}\xspace}
\let\pageoriginaled\pageoriginale

\usepackage{ifthen}
\newboolean{original}
\def\sisi{\ifthenelse{\boolean{original}}}
\setboolean{original}{false}

\newcounter{cnote}
\def\nde#1{\ifthenelse{\boolean{original}}{}{\stepcounter{cnote}\def\thefootnote{\arabic{cnote}}\protect\footnote[\value{cnote}]{N.D.E.: #1}}}
\def\ndemark{\ifthenelse{\boolean{original}}{}{\stepcounter{cnote}\def\thefootnote{\arabic{cnote}}\protect\footnotemark[\value{cnote}]}}
\def\ndetext#1{\ifthenelse{\boolean{original}}{}{\def\thefootnote{\arabic{cnote}}\protect\footnotetext[\value{cnote}]{N.D.E.: #1}}}

\setmainlanguage{french}
\setotherlanguage{english}

% unicode-math deliberately disables the legacy \Arrowvert glyph.  The
% source uses it once as a rotated equality marker; \Vert has the intended
% appearance and keeps the historical diagram buildable with LuaLaTeX.
\renewcommand{\Arrowvert}{\Vert}


%%
%% Pour afficher le numéro de version Git (exige que LuaTeX soit appelé avec --shell-escape)
%%

\catcode`\%=12
\directlua{
function command_output(cmd)
  local output_pipe = io.popen(cmd)
  local line_iter = output_pipe and output_pipe:lines()
  local line = line_iter and line_iter()
  if output_pipe then output_pipe:close() end
  return line
end
local git_command1 = [[git log --pretty=format:"%h %ai" -1]]
local git_command2 = [[if git diff --quiet HEAD ; then echo clean ; else echo dirty ; fi]]
function git_status()
  local status = command_output(git_command1)
  if status then
    tex.tprint({luatexbase.catcodetables.CatcodeTableLaTeX, string.char(0x5c).."texttt{"}, {-2,status}, {luatexbase.catcodetables.CatcodeTableLaTeX, "}"})
    if not (command_output(git_command2) == "clean") then tex.sprint(-2, " (modifications non validées)") end
  else
    tex.sprint(-2, "?")
  end
end
}
\catcode`\%=14
\def\showgitstatus{\directlua{git_status()}}


% macros pour LaTeX
\renewcommand{\mathsf}{}

% Police différente pour U+2113 SCRIPT SMALL L
\setmathfont[range="2113]{latinmodern-math.otf}

% Give every explicit SGA label its own destination and preserve the printed
% number supplied as the optional argument.  The former definition inherited
% whichever counter anchor happened to precede it and ignored #1 entirely.
\makeatletter
\def\etiquette[#1]#2{%
  \phantomsection
  \protected@edef\@currentlabel{#1}%
  \label{#2}%
}
\makeatother
\def\eqhyperref[#1]#2{(\hyperref[#1]{#2})}
\def\pagehyperref[#1]#2{(page \hyperref[#1]{#2})}
\def\hyperrefext[#1][#2]#3{\hyperref[#2]{#1 #3}}
\def\eqhyperrefext[#1][#2]#3{\hyperref[#2]{#1 #3}}     %cf. exposes 17
\def\neweqhyperrefext[#1][#2]#3{\hyperref[#2]{#1 #3}}  %cf. exposes 17

\def\shat#1{\protect\widehat{#1}}
\def\stilde#1{\protect\widetilde{#1}}
%%%%%%%%%%%%%%%%%%


\def\etcn{etc.\xspace}
\def\etco{etc...\xspace}
\def\etc{\sisi{\etco}{\etcn}}
\def\ldotsn{\sisi{...}{\etcn}}

\def\cf{{cf.}\kern.3em}
\def\Cf{{Cf.}\kern.3em}
\def\ie{{i.e.}\kern.3em}
\def\iev{{i.e.},\xspace}
\def\loccito{loc.\kern2pt cit.\xspace}
\def\loccitn{\textit{loc.\kern2pt cit.}\xspace}
\def\loccit{\sisi{\loccito}{\loccitn}}
\def\EGA{\textup{EGA}\kern.2em}
\def\SGA{\textup{SGA}\kern.3em}

\def\quoi{}

\def\quav{\quad,\quad}
\def\quaet{\quad \text{et}\quad}
\def\fibreproduct#1#2#3{#1 ×_{#3} #2}
\let\fprod\fibreproduct
\def\ccheck#1#2{\mbox{${\check{\check{\rm #1}}}^{\lower2pt\hbox{$\scriptstyle #2$}}$}}
\def\sheaf#1{\protect\mathcal{#1}}
\def\vieuxsheaf#1{{#1}} % cf. problème avec sheaf{a\Phi_i} par exemple...
\DeclareMathOperator{\Hom}{Hom}
\DeclareMathOperator{\can}{can}
\DeclareMathOperator{\Ind}{Ind}
\DeclareMathOperator{\ind}{ind}
%\DeclareMathOperator{\Hhomtop}{\mathcal{H}\!\mathit{omtop}}
\DeclareMathOperator{\Hhomtop}{\mathbf{Homtop}}
\DeclareMathOperator{\Point}{Point}
\DeclareMathOperator{\Pro}{Pro}
\DeclareMathOperator{\Top}{Top}
\DeclareMathOperator{\Mod}{Mod}
\DeclareMathOperator{\Cart}{Cart}
%\DeclareMathOperator{\Mmorsite}{\mathcal{M}\!\mathit{orsite}}
\DeclareMathOperator{\Mmorsite}{\mathbf{Morsite}}
\DeclareMathOperator{\loc}{loc}
\DeclareMathOperator{\Loc}{Loc}
\DeclareMathOperator{\Pic}{Pic}
\DeclareMathOperator{\Et}{Et}
\DeclareMathOperator{\card}{card}
\DeclareMathOperator{\Ens}{Ens}
\DeclareMathOperator{\h}{\underline{H}}
\DeclareMathOperator{\Fonct}{Fonct}
\DeclareMathOperator{\Ker}{Ker}
\DeclareMathOperator{\but}{but}
\DeclareMathOperator{\Ob}{Ob}
\DeclareMathOperator{\ob}{ob}
\DeclareMathOperator{\Lib}{Lib}
\DeclareMathOperator{\esj}{esj}
\DeclareMathOperator{\id}{id}
\DeclareMathOperator{\Id}{Id}
\DeclareMathOperator{\ad}{ad}
\DeclareMathOperator{\Iim}{Im}
\DeclareMathOperator{\pr}{pr}
\DeclareMathOperator{\Cat}{Cat}
\DeclareMathOperator{\car}{car}
\DeclareMathOperator{\SysInd}{SysInd}
\DeclareMathOperator{\Sup}{Sup}
\DeclareMathOperator{\Coker}{Coker}
\DeclareMathOperator{\Filt}{Filt}
\DeclareMathOperator{\acc}{acc}
\DeclareMathOperator{\sur}{sur}
\DeclareMathOperator{\cat}{cat}
\DeclareMathOperator{\Ord}{Ord}
\DeclareMathOperator{\Gr}{Gr}
\DeclareMathOperator{\Cov}{Cov}
\DeclareMathOperator{\Caf}{Caf}
\DeclareMathOperator{\coker}{coker}
\DeclareMathOperator{\coim}{coim}
\DeclareMathOperator{\im}{im}
\DeclareMathOperator{\Ouv}{Ouv}
\DeclareMathOperator{\diag}{diag}
\DeclareMathOperator{\Esp}{Esp}
\DeclareMathOperator{\TOP}{TOP}
\DeclareMathOperator{\Unif}{Unif}
\DeclareMathOperator{\Iind}{\mathbf{Ind}}
\DeclareMathOperator{\Hhomcart}{\mathbf{Homcart}}
\DeclareMathOperator{\Hhomtopess}{\mathbf{Homtopess}}
\DeclareMathOperator{\Ppoints}{\mathbf{Points}}
\DeclareMathOperator{\Ppoint}{\mathbf{Point}}
\DeclareMathOperator{\Ttors}{\mathbf{Tors}}
\DeclareMathOperator{\Ssex}{\mathbf{Sex}}
\DeclareMathOperator{\Ffib}{\mathbf{Fib}}
\DeclareMathOperator{\Vvois}{\mathbf{Vois}}
\DeclareMathOperator{\Ppointess}{\mathbf{Pointess}}
\DeclareMathOperator{\scrHom}{\mathbf{Hom}}
\DeclareMathOperator{\sheafPic}{\mathbf{Pic}}
\DeclareMathOperator{\sheafAut}{\mathbf{Aut}}
\DeclareMathOperator{\Hhomtoplocan}{\mathbf{Homtoplocan}}
\DeclareMathOperator{\Hhomtopan}{\mathbf{Homtopan}}
\DeclareMathOperator{\Pptgeom}{\mathbf{Ptgéom}}
\DeclareMathOperator{\Morsite}{Morsite}
\DeclareMathOperator{\Fib}{Fib}
\DeclareMathOperator{\Fais}{Fais}
\DeclareMathOperator{\kar}{kar}
\DeclareMathOperator{\Groupes}{Groupes}
\DeclareMathOperator{\Restr}{Restr}
\DeclareMathOperator{\Isom}{Isom}
\DeclareMathOperator{\Prol}{Prol}
\DeclareMathOperator{\Res}{Res}
\DeclareMathOperator{\Tors}{Tors}
\DeclareMathOperator{\esp}{esp}
\DeclareMathOperator{\es}{es}
\DeclareMathOperator{\od}{od}
\DeclareMathOperator{\gr}{gr}
\DeclareMathOperator{\sob}{sob}
\DeclareMathOperator{\dfn}{dfn}
\DeclareMathOperator{\Sch}{Sch}
\DeclareMathOperator{\ouv}{ouv}
\DeclareMathOperator{\cons}{cons}
\DeclareMathOperator{\Nerf}{Nerf}
\DeclareMathOperator{\Points}{Points}
\DeclareMathOperator{\Espsob}{Espsob}
\DeclareMathOperator{\Homtop}{Homtop}
\DeclareMathOperator{\Pt}{Pt}
\DeclareMathOperator{\Quot}{Quot}
\DeclareMathOperator{\glob}{glob}
\DeclareMathOperator{\Homfib}{Homfib}
\DeclareMathOperator{\Qouv}{Qouv}
\DeclareMathOperator{\Quotouv}{Quotouv}
\DeclareMathOperator{\ens}{ens}
\DeclareMathOperator{\ab}{ab}
\DeclareMathOperator{\an}{an}
\DeclareMathOperator{\fs}{fs}
\DeclareMathOperator{\red}{red}
\DeclareMathOperator{\Ab}{Ab}
\DeclareMathOperator{\cosq}{cosq}
\DeclareMathOperator{\Ann}{Ann}
\DeclareMathOperator{\Zar}{Zar}


%%%%%%%%%%%%
%%%%% SGA4-2
\DeclareMathOperator{\Sex}{Sex}
\DeclareMathOperator{\vois}{vois}
\DeclareMathOperator{\Vois}{Vois}
\DeclareMathOperator{\Ext}{Ext}
\DeclareMathOperator{\ext}{ext}
\DeclareMathOperator{\Fer}{Fer}
\DeclareMathOperator{\Add}{Add}
\DeclareMathOperator{\Cocart}{Cocart}
\DeclareMathOperator{\cosk}{cosk}
\DeclareMathOperator{\plat}{plat}

\DeclareMathOperator{\faisc}{faisc}
\DeclareMathOperator{\sk}{sk}
\DeclareMathOperator{\ch}{ch}

\def\HH{\mathbf{H}}
\def\qc{\mathrm{qc}}
\def\et{\mathrm{et}}
\def\coh{\mathrm{coh}}

\def\Hvr#1{%
\begin{array}[b]{l}
{\kern2pt\scriptstyle r}\\[-8pt]
{\kern1pt\scriptstyle \vee}\\[-5pt]
H^{#1}
\end{array}}

\def\Hvv#1{%
\begin{array}[b]{l}
\kern2pt\rotatebox{90}{${\scriptscriptstyle\ll}$}\\[-6pt]
H^{#1}
\end{array}}

\def\Hvi#1{%
\begin{array}[b]{l}
{\kern0.2pt\scriptscriptstyle \infty}\\[-8pt]
{\kern1pt\scriptstyle \vee}\\[-5pt]
H^{#1}
\end{array}}

\def\Hv#1{%
\begin{array}[b]{l}
{\kern1pt\scriptstyle \vee}\\[-5pt]
H^{#1}
\end{array}}

\DeclareMathOperator{\cart}{cart}
\DeclareMathOperator{\Ffaisc}{\mathbf{Faisc}}


\def\Ev#1{%
\begin{array}[b]{l}
{\kern1pt\scriptstyle \vee}\\[-5pt]
E'_{#1}
\end{array}}

\DeclareMathOperator{\Cont}{Cont}
\DeclareMathOperator{\Ccont}{\mathscr{C}\!\mathit{ont}}
\DeclareMathOperator{\Site}{Site}
\newcommand\MM{\mathbf{M}}
\def\limtop#1{{\displaystyle{\varprojlim_{#1}\!\mbox{\rm top }}}}
\def\Limtop{\mathop{\varprojlim\!\mbox{\rm top}}}
\DeclareMathOperator{\spec}{spec}
\def\ss{\scriptstyle{\mathscr{S}}}
\DeclareMathOperator{\Gal}{Gal}
\def\smallm{\scriptstyle{\mathscr{M}}}
\DeclareMathOperator{\End}{End}
\DeclareMathOperator{\ass}{ass}
\DeclareMathOperator{\zar}{zar}
\DeclareMathOperator{\Ogg}{Ogg}
\DeclareMathOperator{\res}{res}
\DeclareMathOperator{\tr}{tr}
\DeclareMathOperator{\Supp}{Supp}
\DeclareMathOperator{\codim}{codim}
\DeclareMathOperator{\rk}{rk}
\DeclareMathOperator{\cur}{cur}
\DeclareMathOperator{\cdqc}{cdqc}
\newcommand\DD{\mathbf{D}}
\DeclareMathOperator{\cl}{cl}
\def\JJ{\mathbf{J}}
\def\FF{\mathbf{F}}
\def\BB{\mathbf{B}}
\DeclareMathOperator{\Alget}{Alget}
\DeclareMathOperator{\Ass}{Ass}
\DeclareMathOperator{\supp}{supp}
\DeclareMathOperator{\rad}{rad}
\DeclareMathOperator{\prof}{prof}
\DeclareMathOperator{\cint}{int}
\DeclareMathOperator{\sic}{sic}
\DeclareMathOperator{\Tor}{Tor}
\DeclareMathOperator{\tors}{tors}
\DeclareMathOperator{\Sym}{Sym}
\DeclareMathOperator{\parf}{parf}
\DeclareMathOperator{\tor}{tor}
\DeclareMathOperator{\defn}{defn}
\DeclareMathOperator{\Div}{Div}
\DeclareMathOperator{\Pol}{Pol}
\DeclareMathOperator{\Fex}{Fex}
\DeclareMathOperator{\Lie}{Lie}
\DeclareMathOperator{\tot}{tot}
\DeclareMathOperator{\Aut}{Aut}
\DeclareMathOperator{\HOM}{HOM}
\DeclareMathOperator{\pch}{pch}
\DeclareMathOperator{\Ch}{Ch}
\DeclareMathOperator{\EXT}{EXT}
\DeclareMathOperator{\PIC}{PIC}
\DeclareMathOperator{\Rf}{Rf}
\DeclareMathOperator{\inj}{inj}
\newcommand\VV{\mathbf{V}}

\def\Xv{%
\begin{array}[b]{l}
{\kern1pt\scriptstyle \vee}\\[-5pt]
X
\end{array}}

%%%%%%%
\newcommand\TT{\mathbf{T}}
\newcommand\XX{\mathbf{X}}
\newcommand\EE{\mathbf{E}}
%%%%%%

\def\qedsymbol{\textup{C.Q.F.D.}}
\let\CQFD\qed
\let\cqfd\qed

%%%%%\bigdot
\font\bigsymb=cmsy10 at 4pt
\def\bigdot{{\kern1.2pt\raise 1.5pt\hbox{\bigsymb\char15}}}
%%%%%

\def\Lright{\mathop{\vcenter{\hbox{\hbox{$\longrightarrow$}\lower8pt \hbox{\kern-16pt\hbox{$\longrightarrow$}}}}}}

\def\Aa{{\mathcal A}}
\def\Bb{{\mathcal B}}
\def\Cc{{\mathcal C}}
\def\Dd{{\mathcal D}}
\def\Ee{{\mathcal E}}
\def\Ff{{\mathcal F}}
\def\Gg{{\mathcal G}}
\def\Hh{{\mathcal H}}
\def\Ii{{\mathcal I}}
\def\Jj{{\mathcal J}}
\def\Kk{{\mathcal K}}
\def\Ll{{\mathcal L}}
\def\Mm{{\mathcal M}}
\def\Nn{{\mathcal N}}
\def\Oo{{\mathcal O}}
\def\Pp{{\mathcal P}}
\def\Qq{{\mathcal Q}}
\def\Rr{{\mathcal R}}
\def\Ss{{\mathcal S}}
\def\Tt{{\mathcal T}}
\def\Uu{{\mathcal U}}
\def\Vv{{\mathcal V}}
\def\Ww{{\mathcal W}}
\def\Xx{{\mathcal X}}
\def\Yy{{\mathcal Y}}
\def\Zz{{\mathcal Z}}

\newcommand{\CC}{\mathbf{C}}
\newcommand{\GG}{\mathbf{G}}
\newcommand{\KK}{\mathbf{K}}
\newcommand{\LL}{\mathbf{L}}
\newcommand{\NN}{\mathbf{N}}
\newcommand{\PP}{\mathbf{P}}
\newcommand\QQ{\mathbf{Q}}
\newcommand\RR{\mathbf{R}}
\newcommand\ZZ{\mathbf{Z}}

\def\to{\mathchoice{\longrightarrow}{\rightarrow}{\rightarrow}{\rightarrow}}
\def\from{\mathchoice{\longleftarrow}{\leftarrow}{\leftarrow}{\leftarrow}}
\def\mto{\mathchoice{\longmapsto}{\mapsto}{\mapsto}{\mapsto}}
\def\mfrom{\mathrel{\from\joinrel\mapstochar\,}}
\def\hto{\mathrel{\lhook\joinrel\to}}
\def\hfrom{\mathrel{\from\joinrel\rhook}}

\def\implies{\mathchoice{\Longrightarrow}{\Rightarrow}{\Rightarrow}{\Rightarrow}}
\def\implied{\mathchoice{\Longleftarrow}{\Leftarrow}{\Leftarrow}{\Leftarrow}}
\def\iff{\mathchoice{\Longleftrightarrow}{\Leftrightarrow}{\Leftrightarrow}{\Leftrightarrow}}

\def\To#1{\mathchoice{\xrightarrow{\labelstyle~~#1~}}{\stackrel{#1}{\longrightarrow}}{}{}}
\def\Mto#1{\mathrel{\mapstochar\To{#1}}}
\def\From#1{\mathchoice{\xleftarrow{\labelstyle~~#1~}}{\stackrel{#1}{\longleftarrow}}{}{}}

\newcommand{\isomto}{\overset{\sim}{\longrightarrow}}
\newcommand{\isomfrom}{\overset{\sim}{\longleftarrow}}

\newcommand{\IFF}{\unskip~$\iff$~\ignorespaces}
\newcommand{\THEN}{\unskip~$\implies$~\ignorespaces}

\let\leq\leqslant
\let\geq\geqslant
\let\le\leqslant
\let\ge\geqslant

\def\cf{{cf.}\kern.3em}
\def\Cf{{Cf.}\kern.3em}
\def\ie{{i.e.}\kern.3em}
\def\iev{{i.e.},\xspace}
%\def\loccito{loc.\kern2pt cit.\xspace}
%\def\loccitn{\textit{loc.\kern2pt cit.}\xspace}
%\def\loccit{\sisi{\loccito}{\loccitn}}
\def\EGA{\textup{EGA}\kern.2em}
\def\SGA{\textup{SGA}\kern.3em}

\let\oldS\S
\def\S{\oldS\kern .15em}
\def\ptbl{.\kern .2em}
\def\resp{resp.\kern.3em}
\def\Exp{Exp.\kern.3em}

\let\bar\overline
\let\hat\widehat
\let\tilde\widetilde
\let\dpl\displaystyle
\let\ldots\dots
\let\emptyset\varnothing
\let\skpt\skippointrait
\let\epsilon\varepsilon

\def\sfrac#1#2{{#1}/{#2}}

\DeclareMathOperator{\cd}{cd}
%\DeclareMathOperator{\Hhom}{\mathcal{H}\!\mathit{om}}
\DeclareMathOperator{\Hhom}{\mathbf{Hom}}
%\DeclareMathOperator{\Eext}{\mathcal{E}\!\mathit{xt}}
\DeclareMathOperator{\Eext}{\mathbf{Ext}}
\DeclareMathOperator{\Spec}{Spec}
\DeclareMathOperator{\Tr}{Tr}

\let\emptyset=\varnothing

%%%%%%%%%%%%%%%
%%% expose XVII
%%%%%%%%%%%%%%%

\makeatletter
\let\old@optexpose\optexpose
\newcommand{\empty@optexpose}[1]{}
\newcommand{\paren@optexpose}[1]{\old@optexpose{#1}(}
\newcommand{\refnoexpose}[1]{\let\optexpose\empty@optexpose\ref{#1}\let\optexpose\old@optexpose}%\refnoexpose{X.cor5.2} engendre "5.2" au lieu de "X 5.2"
\newcommand{\neweqref}[1]{\let\optexpose\paren@optexpose\textup{\ref{#1})}\let\optexpose\old@optexpose}%\neweqref{XII.eq4.1} engendre "XII (4.1)" au lieu de "(XII 4.1)"
\makeatother
\newcommand{\refvb}[2][VI${}_{\text{B}}$ ]{\sisi{#1\refnoexpose{#2}}{\ref{#2}}}

\newcommand{\leftexp}[2]{{\vphantom{#2}}^{#1}{#2}}
\newcommand{\AAA}{\mathbf{A}}
\newcommand{\Sf}{\mathfrak{S}}
\DeclareMathOperator*{\quotprojlim}{\text{``$\varprojlim$''}}
\DeclareMathOperator*{\quotinjlim}{\text{``$\varinjlim$''}}
\DeclareMathOperator{\Fl}{F\sisi{\ell}{l}}
\DeclareMathOperator{\Pl}{P\sisi{\ell}{l}}
%\DeclareMathOperator{\cl}{c\sisi{\ell}{l}}
\DeclareMathOperator{\tdf}{tdf}
\DeclareMathOperator{\rg}{rg}
\newcommand{\iInd}{\sisi{I}{i}nd}
\newcommand{\iind}{ind}
\newcommand{\ppro}{pro}
\newcommand{\tTor}{\sisi{T}{t}or}
\newcommand{\ttor}{tor}
\newcommand{\chap}{Chap.}
\newcommand{\Ffex}{\Ff\mathit{ex}}
\newcommand{\sheafk}{\underline{k}}
\newcommand{\sheafn}{\underline{n}}
\newcommand{\sheafGamma}{\underline{\Gamma}}

\sisi{
\DeclareMathOperator{\simarrow}{\rule{.15mm}{1.8mm}\kern -3.4pt\rightsquigarrow}
\newcommand{\sheafell}{\underline{\ell}}
\newcommand{\sheafZZ}{\underline{\ZZ}}
}{
\newcommand{\nderef}[1]{la N.D.E.~(\ref{#1}) page~\pageref{#1}}
}

%%%%%%%%%%%%%%%%%
%%%%% expose XVII
%%%%%%%%%%%%%%%%%



\externaldocument{../01/01.tex}
\externaldocument{../02/02.tex}
\externaldocument{../03/03.tex}
\externaldocument{../04/04.tex}
\externaldocument{../05/05.tex}
\externaldocument{../05bis/05bis.tex}
\externaldocument{../06/06.tex}
\externaldocument{../07/07.tex}
\externaldocument{../08/08.tex}
\externaldocument{../09/09.tex}
%\externaldocument{../10/10.tex}
\externaldocument{../11/11.tex}
\externaldocument{../12/12.tex}
\externaldocument{../13/13.tex}
\externaldocument{../14/14.tex}
\externaldocument{../15/15.tex}
\externaldocument{../16/16.tex}
\externaldocument{../17/17.tex}
\externaldocument{../18/18.tex}
\externaldocument{../19/19.tex}

\setboolean{original}{false}

\begin{document}

\begin{center}
XIV. Théorème de finitude pour un morphisme propre ; dimension cohomologique des schémas algébriques affines\\
M. Artin 
\end{center}

\selectlanguage{english}version : \showgitstatus\selectlanguage{french}
\tableofcontents

\else
\chapter{Théorème de finitude pour un morphisme propre ; dimension cohomologique des schémas algébriques affines}
%\addtocontents{toc}{M. Artin}
\begin{center}
M. Artin 
\end{center}
\fi

\setcounter{pageoriginal}{144}

\newcommand{\degtr}{{\rm deg.tr.\ }}    %%% was texed as "\deg. \tr."

%%% The roman numeral in the "table des matieres" should be capital roman letters, as in the original.

%%% Theorem 1.1 is referred to as 1.1, while Lemma 1.4 etc. are referred to by (1.4) etc, 
%%% so with brackets. This is faithful to the original, but typographically I would prefer only one 
%%% of both choices. Sometimes (x.y) is in order to keep a distance from a formula.
%%% To make things worse, also Lemma are sometimes referred to as x.y instead of (x.y).

%%% The citattion EGA x.y or other SGA? should also turned into hyperlinks!

%%% Etiquette: There is a \ref{XVII} which has no \label

Cet\pageoriginale exposé contient deux théorèmes importants qui se
démontrent en 
utilisant le théorème de changement de base pour un morphisme propre.

\section{Théorème de finitude pour un morphisme propre}\etiquette[1]{XIV.sec1}

\begin{enonce*}{Théorème 1.1}\etiquette[1.1]{XIV.thm1.1}
Soit $f:X\rightarrow S$ un morphisme propre et de présentation
finie. Soit $A$ un anneau noethérien à gauche qui est une
$\ZZ/n$-algèbre pour $n\in \NN$ convenable. Soit $F$ un
faisceau d'ensembles (\resp de groupes, \resp de $A$-modules)
\emph{constructible} sur $X$. Alors le faisceau $f_{\ast}F$ (\resp
$R^1f_{\ast}F$, \resp $R^qf_{\ast}F$ pour chaque $q\geq 0$) 
%%% typographical remark: 
%%% The brackets in the \enonce{} are badly tilted and do not keep the right spacing to the text.
est également constructible.
\end{enonce*}

En particulier, on a:

\begin{enonce*}{Corollaire 1.2}\etiquette[1.2]{XIV.cor1.2}
Soit $X$ un schéma propre sur le spectre d'un corps séparablement clos
$k$, et soit $F$ un faisceau abélien de torsion constructible sur
$X$. Alors les groupes $H^q(X,F)$ sont finis pour chaque $q\geq 0$.
\end{enonce*}

\begin{enonce*}[remark]{Remarque 1.3}\etiquette[1.3]{XIV.rem1.3}
Le théorème est faux pour les faisceaux abéliens si l'on omet la
condition que $f$ soit propre, comme on voit en prenant
$F=\ZZ/p$, $p=\car k$, $X=\EE^1_k$, l'espace affine, et
$f$ l'inclusion de $X$ dans $\PP^1_k$, ou le morphisme
structural $X\rightarrow \Spec k$. Cependant, on 
conjecture\nde{Ofer Gabber a présenté  la généralisation suivante de \ref{XIV.thm1.1} lors des conférences en l'honneur de Luc Illusie et de Pierre Deligne en 2005. Soit $f:X \to Y$ un morphisme de type fini de schémas noethériens quasi-excellents (\EGA IV 7.8), et soit $F$ un faisceau constructible de $\mathbf{Z}/n$-modules sur $X_{\text{ét}}$ avec $n$ inversible sur $Y$. Alors, les faisceaux $R^i f_\ast F$ sur $Y_{\text{ét}}$ sont constructibles et nuls sauf pour un nombre fini de $i$. Il y a des assertions analogues pour faisceaux constructibles d'ensembles \resp de groupes d'ordre inversible sur $Y$. \newline
Avant le résultat de Gabber, le théorème \ref{XIV.thm1.1} était connu dans le cas suivant. Soit $S$ un schéma régulier de dimension $\leq1$ et soit $f: X \to Y$  un morphisme de $S$-schémas de type fini et $F$ un faisceau de $A$-modules à gauche sur $X_{\text{ét}}$. Alors, les faisceaux $R^i f_\ast F$ sont constructibles (\SGA 4 1/2 Finitude 1.1).}
qu'il est
vrai si $F$ est premier aux\pageoriginale caractéristiques résiduelles
de $S$ et $S$ 
est « excellent » \EGA IV 7.8. On peut le démontrer, lorsqu'on
dispose du théorème de résolution des singularités (par exemple en
caractéristique nulle) pour un schéma excellent $S$ d'égales
caractéristiques (\hyperrefext[XIX][XIX.thm5.1]{5.1}).
\end{enonce*}

\noindent
\emph{Démonstration de} \hyperref[XIV.thm1.1]{1.1}. L'assertion est locale sur
$S$ et on peut 
donc prendre $S$ affine. Alors $S$ est limite des spectres d'anneaux
de type fini sur $\ZZ$, et les données du théorème sont telles
qu'on peut les obtenir, par changement de base $S\rightarrow
S_{0}$, d'un morphisme propre $f_{0}:X_{0}\rightarrow
S_{0}$ avec $S_{0}$ de type fini sur $\ZZ$ et un
faisceau constructible $F_{0}$ sur $X_{0}$, \cf \EGA IV 8 et 
\hyperrefext[IX][IX.cor2.7.4]{2.7.4}. Or, puisque l'image inverse d'un faisceau constructible est
encore constructible (\hyperrefext[IX][IX.prop2.4]{2.4} (iii)), on est ramené, en appliquant 
\hyperrefext[XII][XII.thm5.1]{5.1} au morphisme de changement de base $S_{0}\leftarrow S$, à
démontrer la constructibilité de $f_{0 \sisi{}{,} \ast} F_{0}$ (\resp de
$R^{q}f_{0 \sisi{}{,} \ast} F_{0}$), d'où le

\begin{enonce*}{Lemme 1.4}\etiquette[1.4]{XIV.lem1.4}
Il suffit de vérifier \hyperref[XIV.thm1.1]{1.1} avec $S$ de type fini sur $\Spec \ZZ$.
\end{enonce*}

\begin{enonce*}{Lemme 1.5}\etiquette[1.5]{XIV.lem1.5}
L'assertion ensembliste de \hyperref[XIV.thm1.1]{1.1} est vraie, c'est-à-dire $f_{\ast}F$
est un faisceau constructible d'ensembles si $F$ l'est.
\end{enonce*}

\noindent
\emph{Démonstration}. On peut supposer $S$ noethérien (\hyperref[XIV.lem1.4]{1.4}). Soit
$F\rightarrow G=\prod_i\pi_{i \sisi{}{,} \ast}(C_i)$ une injection avec
$\pi_i:X_i\rightarrow X$ fini et $C_i$ constant et constructible sur
$X_i$ (\hyperrefext[IX][IX.prop2.14]{2.14}). On a $f_{\ast}F\subset f_{\ast}G$ 
%%% new line break within formula in text: spend extra line to induce new line breaks?
et par suite 
(\hyperrefext[IX][IX.prop2.9]{2.9} \sisi{(iii)}{(ii)}) il suffit de traiter le cas $F=G$. Or
$f_{\ast}G=\prod_i 
(f\pi_i)_{\ast}G_i$, et un produit fini de faisceaux constructibles
sur $S$ est constructible (\sisi{\hyperrefext[IX][IX.prop2.14]{2.14}}{\hyperrefext[IX][IX.prop2.6]{2.6} (ibis)}). En remplaçant $X$ par $X_i$,
on est ramené au cas où $F=D_X$ est un faisceau constant
constructible, à valeur $D$. Or en vertu de \hyperrefext[XII][XII.thm5.1]{5.1} (i), la fibre
$f_{\ast}(F)_{\bar{s}}$ est isomorphe à $D^{C(s)}$, où
$C(s)=\pi_0(X_{\bar{s}})$,\pageoriginale donc elle est
finie. D'après \hyperrefext[IX][IX.prop2.13]{2.13} (iii) il suffit de démontrer que la fonction
« nombre d'éléments de la fibre de $f_{\ast}F$ »  est
constructible sur $S$. Or cette fonction est $\card(D^{C})$, où $C$
est la fonction « nombre de composantes connexes dans la fibre
géométrique », et $C$ est constructible (\EGA IV 9.7.9), d'où le
résultat.

\begin{enonce*}{Lemme 1.6}\etiquette[1.6]{XIV.lem1.6}
Soit $S$ un schéma noethérien, $f:X\rightarrow S$ propre, $F$ un
faisceau de groupes (\resp de $A$-modules) constructible sur $X$, $q$
un entier égal à $0$ ou $1$ (\resp un entier $\geq 0$). Soient $S_i$
$(1\leq i\leq n)$ des sous-schémas de $S$ tels que $S=\cup_i S_i$ et
soient $F_i$, $f_i:X_i\rightarrow S_i$ les « restrictions »  des
données aux $S_i$. Alors si pour tout $i$, $R^{q}f_{i \sisi{}{,} \ast}F_i$ est
constructible, il en est de même de $R^{q}f_{\ast}F$. 
\end{enonce*}

\noindent
\emph{Démonstration}. La restriction de $R^{q}f_{\ast}F$ à $S_i$ est
isomorphe à $R^{q}f_{i \sisi{}{,} \ast}F_i$ d'après \hyperrefext[XII][XII.thm5.1]{5.1}, et le lemme suit
de \hyperrefext[IX][IX.prop2.8]{2.8}.

On procède maintenant par des raisonnements analogues à ceux de \hyperrefext[XII][XII.sec8]{8}:

\begin{enonce*}{Lemme 1.7}\etiquette[1.7]{XIV.lem1.7}
Soient $S$ noethérien,
$$
\xymatrix@=1.25cm{
X        \ar[rd]_{f} &           &\bar{X} \ar[ll]_{\pi}
 \ar[ld]^{\bar{f}}\\
                 &    S      &         
}
$$
un diagramme commutatif de morphismes propres, et $j:U\rightarrow X$
un ouvert tel que l'ouvert $\bar{U}=\fprod{U}{\bar{X}}{X}$ s'envoie
isomorphiquement sur $U$. Soit $Y=X-U$ le schéma fermé réduit et
soient $f_{0}:Y\rightarrow S$ et $i:Y\rightarrow X$ les morphismes
canoniques. Supposons le théorème \hyperref[XIV.thm1.1]{1.1} 
vrai\nde{L'hypothèse sur $R^q \pi_\ast$ ne figure que pour $q=0$ dans la 
démonstration. Dans ce cas, on sait déjà la constructibilité par (\hyperref[XIV.lem1.5]{1.5}).}
pour $\bar{f}$, $f_{0}$ et $\pi$. Alors il est également vrai pour \sisi{}{$f$.}
\end{enonce*}
%%% In Lemma 1.7, in my opinion the assumption on pi is never used 
%%% (except for q=0 which is proved already in Lemma 1.5)

\noindent
\emph{Démonstration}.\pageoriginale Traitons d'abord le cas d'un
faisceau de $A$-modules $F$ sur $X$. On a la suite exacte
$$
0\longrightarrow j_{!}j^{\ast}F\longrightarrow F\longrightarrow
i_{\ast}i^{\ast}F\longrightarrow 0{\quoi},
$$
et compte tenu de la suite exacte de cohomologie correspondante pour
les $R^{q}f_{\ast}$, et de \hyperrefext[IX][IX.prop2.6]{2.6}, on voit qu'il suffit de
vérifier la 
constructibilité des $R^{q}f_{\ast}$ pour les membres extrêmes (notons
que $i_{\ast}i^{\ast}F$ est constructible d'après \hyperrefext[IX][IX.prop2.4]{2.4} (iii) et 
\sisi{IX 2.15}{\hyperrefext[IX][IX.prop2.14]{2.14}} (i), donc $j_{!}j^{\ast}F$ l'est aussi
(\hyperrefext[IX][IX.prop2.6]{2.6}). Puisque $i$ est 
une immersion fermée,
$$
(R^{q}f_{\ast})i_{\ast}(i^{\ast}F)\simeq
(R^{q}f_{0 \sisi{}{,} \ast})(i^{\ast}F){\quoi},
$$
(\hyperrefext[VIII][VIII.cor5.6]{5.6}), d'où la constructibilité dans ce cas d'après l'hypothèse
sur $f_{0}$. De plus, désignant par
$\bar{j}:\bar{U}\rightarrow \bar{X}$ le morphisme d'inclusion et
$\bar{F}=\pi^{\ast}F$, on a
\begin{equation*}
j_{!}j^{\ast}F\simeq \pi_{\ast}\bar{j}_{!}\bar{j}^{\ast}\bar{F}{\quoi},\tag{$\ast$}
\end{equation*}
comme on vérifie en définissant d'abord une flèche $\rightarrow$, et
prouvant qu'elle est inversible fibre par fibre en utilisant \hyperrefext[XII][XII.thm5.1]{5.1}\sisi{~(i).}{.} En chaque point géométrique $\bar{y}$ de $Y$, la restriction de
$\bar{j}_{!}\bar{j}^{\ast}\bar{F}$ à la fibre $\bar{X}_{\bar{y}}$ de
$\bar{X}$ sur $X$ est nulle, et par suite \hyperrefext[XII][XII.thm5.1]{5.1} \sisi{(iii)}{} la fibre de
$(R^{q}\pi_{\ast})$ $(\bar{j}_{!}\bar{j}^{\ast}\bar{F})$ en $\bar{y}$
est nulle pour chaque $q\geq 0$. Puisque d'autre part $\pi$ induit un
isomorphisme au-dessus de $X-Y$, on a $(R^{q}\pi_{\ast})$
$(\bar{j}_{!}\bar{j}^{\ast}\bar{F})=0$ si $q>0$. De la suite spectrale
de Leray
$$
E^{p,q}_2=(R^{p}f_{\ast})(R^{q}\pi_{\ast})(\bar{j}_{!}\bar{j}^{\ast}\bar{F})\Longrightarrow
R^{n}\bar{f}_{\ast}(\bar{j}_{!}\bar{j}^{\ast}\bar{F})
$$
et $(\ast)$ on déduit des isomorphismes
$$
(R^{q}\bar{f}_{\ast})(\bar{j}_{!}\bar{j}^{\ast}\bar{F})\simeq
(R^{q}f_{\ast})(j_{!}j^{\ast}F)
$$
pour $q\geq 0$. Puisque $\bar{j}_{!}\bar{j}^{\ast}\bar{F}$ est
constructible (même raisonnement que pour 
\sisi{$\bar{j}_{!}\bar{j}^{\ast}\bar{F}$}{$j_{!}j^{\ast} F$}) le membre de 
\sisi{droite}{gauche}
 l'est aussi,
d'où la constructibilité de
$(R^{q}f_{\ast})(j_{!}j^{\ast}F)$\pageoriginale et donc de
$R^{q}f_{\ast}F$.

Supposons maintenant que $F$ soit un faisceau de groupes constructible
sur $X$. Remplaçons $\pi:\bar{X}\rightarrow X$ par $\pi\amalg i:
(\bar{X}\amalg Y)\rightarrow X$.
%%% typographical remark: I find \amalg too big as a symbol (but I don't know a better symbol).
%%% Better spacing should be ensured at least.
 Il suffit ainsi de démontrer que si
$\pi:\bar{X}\rightarrow X$ est \emph{surjectif} et si 1.1 est vrai
pour $\pi$ et pour $\bar{f}$, il l'est également pour $f$. Soit
d'abord $\bar{F}$ un faisceau de groupes constructible sur
$\bar{X}$. Alors $\pi_{\ast}\bar{F}$ est constructible sur $X$
(\hyperref[XIV.lem1.5]{1.5}) 
et on a une injection $(R^{1}f_{\ast})\pi_{\ast}\bar{F}\hookrightarrow
(R^{1}\bar{f}_{\ast})\bar{F}$. Le membre de droite est constructible
par hypothèse, et par suite le membre de gauche l'est aussi (\hyperrefext[IX][IX.prop2.9]{2.9}
(ii)). Or\sisi{}{,} puisque $\pi$ est surjectif, on a une injection
$F\rightarrow \pi_{\ast}\pi^{\ast}F=G$. Le faisceau
$\bar{F}=\pi^{\ast}F$ est constructible, donc $G$ et $R^{1}f_{\ast}G$
sont aussi constructibles, comme on a vu ci-dessus. On est ainsi
ramené au lemme suivant:

\begin{enonce*}{Lemme 1.8}\etiquette[1.8]{XIV.lem1.8}
Soient $S$ noethérien, $f:X\rightarrow S$ un morphisme propre, et 
$\displaystyle{F\mathop{\hookrightarrow}^{u}G}$
une injection de faisceaux de groupes constructibles sur $X$. Si
$R^{1}f_{\ast}G$ est constructible, $R^{1}f_{\ast}F$ l'est aussi.
\end{enonce*}

\noindent
\emph{Démonstration}. Soit $C=G/F$ qui est un faisceau constructible
d'espaces homogènes sous $G$ (appliquer \hyperrefext[IX][IX.prop2.6]{2.6}). Considérons la suite
exacte
\begin{equation*}
f_{\ast}C\xrightarrow{\delta}R^{1}f_{\ast}F\xrightarrow{u^{1}}R^{1}f_{\ast}G{\quoi}.\tag{$\ast$}
\end{equation*}
Pour démontrer que $R^{1}f_{\ast}F$ est constructible, il suffit par
récurrence noethérienne et \hyperref[XIV.lem1.6]{1.6} de démontrer que si 
$S\neq ∅$, il 
existe un ouvert non-vide de $S$ sur lequel $R^{1}f_{\ast}(F)$ soit
constructible. On peut donc supposer $R^{1}f_{\ast}G$ localement
constant, donc puisque la constructibilité est une notion locale pour
la topologie étale, on peut supposer que $S$ est irréductible et
$R^{1}f_{\ast}G$ est constant (et constructible), à valeur
$E=\{e_1,\ldots,e_n\}$. Soit $D_i$ $(i=1,\ldots,n)$ le
sous-faisceau\pageoriginale de $R^{1}f_{\ast}F$ image inverse de la
section $e_{i}$ par $u^{1}$. On a, puisque dans un topos « les
sommes sont universelles », $R^{1}f_{\ast}F=\amalg_i D_i$, et il
suffit donc de démontrer que chaque $D_i$ est constructible sur un
$S'$ convenable étale sur $S$ et non-vide.

Rappelons que $R^{1}f_{\ast}F$ est le faisceau associé au préfaisceau
%%% new line break within formula in text: spend extra line to induce new line breaks?
$\Rr^{1}F$, où $\Rr^{1}F(S')=H^{1}(\sisi{X'}{X \times_S S'},F)$.
%%% The symbol X' is not defined in the text, although its meaning is determined by 
%%% "functoriality of notation". Nevertheless, I suggest to replace X' by X \times_S S' here.
Par suite
on peut supposer (en remplaçant $S$ par un $S'$ non-vide étale sur
$S$) que $D_i$ est le faisceau « vide »  (donc constructible),
\emph{ou} qu'il existe un élément $\alpha\in H^{1}$ $(X,F)$ qui induit
une section $\sisi{\alpha}{\bar{\alpha}}$ de $D_{i}$. Dans ce dernier cas, soit
$$
f_{\ast}C^{\alpha}\xrightarrow{\delta}R^{1}f_{\ast}F^{\alpha}\longrightarrow
R^{1}f_{\ast}G^{\alpha}
$$
la suite exacte déduite de $(\ast)$ « en tordant à l'aide de
$\alpha$ ». On a
$$
R^{1}f_{\ast}F^{\alpha}\simeq R^{1}f_{\ast}F\quav
R^{1}f_{\ast}G^{\alpha}\simeq R^{1}f_{\ast}G{\quoi},
$$
et la section de $R^{1}f_{\ast}F^{\alpha}$ correspondant à
$\bar{\alpha}$ est la section unité. Donc le sous-faisceau
$D_{i}^{\alpha}$ de $R^{1}f_{\ast}F^{\alpha}$ correspondant à $D_i$
est l'image inverse de la section unité de $R^{1}f_{\ast}G^{\alpha}$,
\ie, $D^{\alpha}_i$ est l'image de $f_{\ast}G^{\alpha}$ par
$\delta$. Puisque $f_{\ast}G$ est constructible (\hyperref[XIV.lem1.5]{1.5}) il en est de
même de $D^{\alpha}_i$, donc de $D_{i}$ (appliquer \hyperrefext[IX][IX.prop2.6]{2.6}), d'où le
lemme.

\begin{enonce*}{Lemme 1.9}\etiquette[1.9]{XIV.lem1.9}
Soit $S$ noethérien et
$$
\xymatrix@=1.25cm{
X       \ar[rd]_{f}      &      &\bar{X} \ar[ll]_{\pi} \ar[ld]^{\bar{f}}\\
              & S    &         
}
$$
un diagramme commutatif de morphismes propres. Si le théorème \hyperref[XIV.thm1.1]{1.1}
est vrai pour\pageoriginale $f$ et $\pi$, il l'est également pour $\bar{f}$.
\end{enonce*}

\noindent
\emph{Démonstration}. Dans le cas d'un faisceau $\bar{F}$ de
$A$-modules constructible, on utilise la suite spectrale de Leray
$$
E^{p,q}_2=(R^{p}f_{\ast})(R^{q}\pi_{\ast}) \bar{F}\Longrightarrow 
(R^{n}\bar{f}_{\ast})\bar{F}{\quoi}.
$$
Par hypothèse, on trouve que $E^{p,q}_2$ est constructible pour chaque
$p$, $q$, et par suite l'aboutissement l'est aussi, comme il résulte
alors de \hyperrefext[IX][IX.prop2.6]{2.6}.

Soit $\bar{F}$ un faisceau de groupes constructible. On la suite
exacte (\hyperrefext[XII][XII.prop3.2]{3.2})
\begin{equation*}
0\longrightarrow (R^{1}f_{\ast})\pi_{\ast}\bar{F}\longrightarrow
(R^{1}\bar{f}_{\ast})\bar{F}\longrightarrow f_{\ast}(R^{1}\pi_{\ast})\bar{F}{\quoi},\tag{$\ast$}
\end{equation*}
où les membres extrêmes sont constructibles d'après l'hypothèse et
\hyperref[XIV.lem1.5]{1.5}. On procède comme dans la démonstration de \hyperref[XIV.lem1.8]{1.8}: Il suffit par
récurrence noethérienne et \hyperref[XIV.lem1.6]{1.6} de supposer $S\neq ∅$, et de
démontrer la constructibilité sur un ouvert non-vide convenable. On
peut donc supposer $f_{\ast}(R^{1}\pi_{\ast})\bar{F}$ localement
constant, et (par localisation étale) même constant et constructible
(\sisi{\hyperrefext[IX][IX.prop2.8]{2.8}}{\hyperrefext[IX][IX.prop2.4]{2.4}~(i)}) à valeur $E=\{e_1,\ldots,e_n\}$. Soit $D_i$ le sous-faisceau
de $(R^{1}\bar{f}_{\ast})\bar{F}$ image inverse de $e_i$, de sorte que
$(R^{1}\bar{f}_{\ast})$ $\bar{F}=\coprod_{i}D_i$. Il suffit de
démontrer que chacun des faisceaux $D_i$ induit un faisceau
constructible sur un $S'/S$ étale connexe et non-vide convenable. Si
$D_i$ n'est pas le « faisceau vide », on se ramène au cas où il
existe un élément $\alpha \in H^{1}$ $(\bar{X},\bar{F})$ qui induit
une section $\bar{\alpha}$ de $D_i$. Soit
$$
\sisi{
0\longrightarrow (R^{1}f_{\ast})\pi_{\ast}\bar{F}\longrightarrow
(R^{1}\bar{f}_{\ast})\bar{F}\longrightarrow
f_{\ast}(R^{1}\pi_{\ast})\bar{F} 
}
{0\longrightarrow (R^{1}f_{\ast})\pi_{\ast}\bar{F}^\alpha \longrightarrow
(R^{1}\bar{f}_{\ast})\bar{F}^\alpha \longrightarrow
f_{\ast}(R^{1}\pi_{\ast})\bar{F}^\alpha}
$$
la suite exacte déduite de $(\ast)$ en tordant $\bar{F}$ à l'aide de
$\alpha$. La section de $(R^{1}\bar{f}_{\ast})$ $\bar{F}^{\alpha}$
correspondant à $\bar{\alpha}$ est la section unité, et par suite le
sous-faisceau $D^{\alpha}_i$ de \hbox{$(R^{1}\bar{f}_{\ast})
\bar{F}^{\alpha}$} correspondant à $D_i$ est l'image inverse de la
section unité de $f_{\ast}(R^{1}\pi_{\ast})$
$\bar{F}^{\alpha}$,\pageoriginale d'où $D^{\alpha}_i\simeq
(R^{1}f_{\ast}) \pi_{\ast}\bar{F}^{\alpha}$, 
%%% new line break within formula in text: spend extra line to induce new line breaks?
qui est bien un faisceau constructible, \cqfd
%%% Typographical remark: In the original, the cqfd is in small letters. That looks much nicer.

\begin{enonce*}{Lemme 1.10}\etiquette[1.10]{XIV.lem1.10}
Il suffit de vérifier \hyperref[XIV.thm1.1]{1.1} dans le cas où $S$ est de type fini sur
$\Spec \ZZ$, et $f$ est projectif et de dimension relative $\leq
1$.
\end{enonce*}

\noindent
\emph{Démonstration}. Supposons d'abord que $f$ soit projectif et
procédons par récurrence sur la dimension relative de $f$, qu'on peut
supposer finie, quitte à remplacer $S$ par un ouvert affine. Si $f$
est de dimension relative $\leq n$, avec $n\geq 2$, on peut trouver,
localement sur $S$, un morphisme fini $\pi:X\rightarrow
\PP^n_S$, et on se ramène par \hyperref[XIV.lem1.9]{1.9} et \sisi{IX 2.15}{\hyperrefext[IX][IX.prop2.14]{2.14}} (i) au cas où $f$
est le morphisme structural de $\PP^n_S$. Considérons le
diagramme commutatif de morphismes
$$
\xymatrix@C=2cm{
                   & \bar{X}\ar[ld]_{\sisi{}{\pi}} \ar[rd]^{a} \ar@{-}[dd]^{\bar{f}} &       \\
\PP^{n}_S \ar[rd]_{f}  &                  &\PP^{1}_S \ar[ld]^{b}\\
                   &  S                          &    
}
$$
déjà envisagé dans la démonstration de \hyperrefext[XII][XII.lem8.3]{8.3}, où $\bar{X}$ est
déduit de $\PP^n_S$ en faisant éclater le sous-schéma fermé
$Y\simeq \PP^{n-2}_S$. La dimension relative de $b$ est $\leq
1$, celle de $a$ est $\leq n-1$, et par suite le théorème est vrai pour
$a$ et $b$ par l'hypothèse de récurrence, donc pour $\bar{f}$ par
(\hyperref[XIV.lem1.9]{1.9}). De plus, la dimension de $\pi$ 
%%% This \pi is different from the \pi used in thebeginning of the proof. Maybe some confusion 
%%% can be avoided by using a different name for one of them (the first one, I suppose).
est $\leq 1$, celle de
$f_{0}:Y\rightarrow S$ est $\leq n-1$. Il résulte donc de
\hyperref[XIV.lem1.7]{1.7} que 
\hyperref[XIV.thm1.1]{1.1} est vrai pour $f$, donc pour tout morphisme projectif. 
Soit
maintenant $f$ propre arbitraire, et procédons par récurrence
noethérienne sur $X$: on peut supposer le théorème vrai pour chaque
sous-schéma fermé distinct de $X$ et que $X\neq ∅$. Soit
$$
\xymatrix@=1.5cm{
X  \ar[rd]_{f}          &           &\bar{X} \ar[ll]_{\pi} \ar[ld]^{\bar{f}}\\
                        &    S       &       
}
$$
un\pageoriginale diagramme du type donné par le lemme de Chow 
\hyperrefext[XII][XII.lem7.1]{7.1}. En appliquant \hyperref[XIV.lem1.7]{1.7} à ce diagramme, 
on déduit \hyperref[XIV.thm1.1]{1.1} pour $f$, d'où le lemme.

\setcounter{subsection}{10}
\subsection{}\etiquette[1.11]{XIV.subsec1.11}
Nous pouvons maintenant achever la démonstration de \hyperref[XIV.thm1.1]{1.1}. Si $F$ est un
faisceau de groupes constructibles, on peut trouver une injection
$F\hookrightarrow G=\prod_i\pi_{i \sisi{}{,} \ast}C_i$ avec $\pi_{i}$ fini et
$C_i$ constant et constructible (\sisi{IX 2.15}{\hyperrefext[IX][IX.prop2.14]{2.14}}), et il résulte de \hyperref[XIV.lem1.8]{1.8} qu'on
peut supposer $F$ constant. De même, si $F$ est un faisceau de
$A$-modules constructible, on peut, en appliquant \sisi{IX 2.15}{\hyperrefext[IX][IX.prop2.14]{2.14}}, trouver une
résolution $F\rightarrow G^{\bigdot}$, $G^{\bigdot}=\{0\rightarrow
G^{0}\rightarrow G^{1}\rightarrow G^{2}\rightarrow\ldots\}$, 
%%% new line break within formula in text: spend extra line to induce new line breaks?
de $F$, telle que chaque $G^{\nu}$ soit de la forme
$G^{\nu}=\prod_i\sisi{\eta}{\pi}_{i \sisi{}{,} \ast}C_i$, 
%%% The original $\eta$ is never defined or specified (these should be finite maps) 
%%% while the $\pi_i$ have been used in the analogous argument above.
%%% However, I would even suggest using $p_i$ as in Proposition IX 2.14(ii) that is 
%%% referred to. Anyway, those maps do not enter again and so the decision of
%%% how to name them is marginal.
et on se réduit par la suite
spectrale de la résolution $G^{\bigdot}$,
$$
E^{p,q}_2=H^{\sisi{q}{p}}((R^{q}f_{\ast})G^{\bigdot})\Longrightarrow
(R^{n}f_{\ast})F{\quoi},
$$
au cas où $F$ est constant et constructible, à valeur $M$.

Rappelons que $A$ est une $\ZZ/n$-algèbre. Soit $\ell$ un
nombre premier qui divise $n$. On a la suite exacte
$$
0\longrightarrow \ell M\longrightarrow M\longrightarrow M/\ell
M\longrightarrow 0{\quoi},
$$
et il suffit encore de prouver \hyperref[XIV.thm1.1]{1.1} pour $\ell M$ et $M/\ell M$.

On se réduit alors, par une récurrence facile, au cas $A$ est une
$\ZZ/\ell$-algèbre, $\ell \in \PP$, et où $M$ est un
$A$-module de type fini. Alors $M$ est un espace vectoriel sur
$\ZZ/\ell$, et par suite le morphisme canonique
$$
H^{q}(Y,\ZZ/\ell)\otimes_{\ZZ/\ell}M\longrightarrow
H^{q}(Y,M)
$$
est bijectif, quel que soit $Y$ quasi-compact et quasi-séparé (en
effet, c'est trivial si $M$ est de rang fini, et les deux membres
commutent aux limites inductives\pageoriginale 
(\sisi{VI 1.6}{\hyperrefext[VI][VI.thm5.1]{5.1} et \hyperrefext[VII][VII.prop3.3]{3.3}})), donc
$$
[R^{q}f_{\ast}(\ZZ/\ell)]\otimes_{\ZZ/\ell}M\xrightarrow{\sim}
R^{q}f_{\ast}(M){\quoi}.
$$
Il suffit donc évidemment de démontrer que $R^{q}f_{\ast}(\ZZ/\ell)$
est constructible en tant que
$\ZZ/\ell$-Module, c'est-à-dire, on est réduit au cas
$A=\ZZ/\ell$, ce que nous supposons désormais.

Notons que le cas où $F$ est de dimension relative $\leq 0$ (\ie où
$f$ est fini) et $F$ est arbitraire résulte maintenant de \hyperrefext[VIII][VIII.cor5.6]{5.6} et
de \hyperref[XIV.lem1.5]{1.5}.

Procédons par récurrence noethérienne sur $S$. D'après \hyperref[XIV.lem1.6]{1.6}, on peut
supposer $S\neq ∅$, et se permettre de remplacer $S$ par un ouvert
non-vide quelconque. De plus, on peut supposer $S$ intègre, de point
générique $s$. La fibre $X_s$ est un schéma algébrique de dimension
$\leq 1$, et par suite il existe une extension radicielle $K'$ de
$k(s)=K$, telle que le normalisé $Y$ de $(X_s\otimes_K K')_{\text{\rm
réd}}$ soit lisse au-dessus de $\Spec K'$ \footnote[1]{\Cf \EGA IV
17.15.14.}. En remplaçant $S$ par un ouvert non-vide, on peut
supposer qu'il existe un morphisme fini radiciel surjectif
$S'\rightarrow S$ tel que $S'_s$ soit $K$-isomorphe à $\Spec K'$, et
un morphisme fini surjectif, un morphisme lisse et projectif
$\bar{X}'\rightarrow S'$ et un $S'$-morphisme fini surjectif
$\pi:\bar{X}'\rightarrow X'=\fprod{X}{S'}{S}$ 
%%% new line break within formula in text: spend extra line to induce new line breaks?
qui induise
$Y\rightarrow X_s \otimes_K K'$ sur les fibres en le point générique
$s$ de $S$ (\EGA IV 9.6.1). Or\sisi{}{,} d'après \hyperrefext[VIII][VIII.thm1.1]{1.1}\sisi{}{,} il est inoffensif de
faire une extension radicielle. On peut donc remplacer $S$ par $S'$ et
$X$ par $X'$, c'est-à-dire, on peut supposer que $f$ s'insère dans un
diagramme
$$
\xymatrix@=1.5cm{
X   \ar[rd]_{f}& &\bar{X} \ar[ll]_{\pi} \ar[ld]^{\bar{f}}\\
         &         S        &           ,
}
$$
où\pageoriginale $\pi$ est surjectif et fini, et $\bar{f}$ est
lisse. De plus, d'après la construction, on peut supposer qu'il existe
un sous-schéma fermé $Y$ de $X$ tel que la fibre générique de $Y$ soit
finie, et tel que $\pi$ induise un isomorphisme au-dessus de $X-Y$. En
remplaçant encore $S$ par un ouvert non-vide, on se ramène au cas
où de plus $f_{0}:Y\rightarrow S$ est un morphisme fini.

Comme $F$ est constant, $\pi^{\ast}F$ est également constant, et
$F\rightarrow \pi_{\ast}\pi^{\ast}F$ est injectif parce que $\pi$ est
surjectif. Si $F$ est un faisceau de groupes, il suffit de démontrer
que $(R^{1}f_{\ast})(\pi_{\ast}\pi^{\ast}F)\simeq
(R^{1}\bar{f}_{\ast})(\pi^{\ast}F)$ est constructible (\hyperref[XIV.lem1.8]{1.8}). Si $F$
est un faisceau de $\ZZ/\ell$-modules, on a une suite exacte
$$
0\longrightarrow F\longrightarrow \pi_{\ast}\pi^{\ast}F\longrightarrow
C\longrightarrow 0{\quoi},
$$
où $C$ est concentré sur $Y$, parce que $\pi$ induit un isomorphisme
au-dessus de $X-Y$. Il s'ensuit que
$R^{q}f_{\ast}C=R^{q}f_{0 \sisi{}{,} \ast}C$ est constructible (nul si
$q>0$), et ainsi il suffit encore de démontrer que $(R^{q}f_{\ast})$
$(\pi_{\ast}\pi^{\ast}F)=(R^{q}\bar{f}_{\ast})$ $(\pi^{\ast}F)$ est
constructible. En remplaçant $f$ par $\bar{f}$ et $F$ par
$\pi^{\ast}F$, on est réduit au cas $F$ constant et constructible et
$f$ est lisse.

Traitons le cas d'un faisceau de groupes constant, $\sisi{f}{F}=G_X$. Toutes les
fibres de $(R^{1}f_{\ast})$ $G_X$ sont finies; en effet, en appliquant
\hyperrefext[XII][XII.cor5.2]{5.2} (ii) on est réduit, pour le prouver, au cas où $S$ est le
spectre d'un corps séparablement clos et où $X$ est lisse et de
dimension $\leq 1$ au-dessus de $S$. Le résultat est connu dans ce
cas (\SGA 1 X 2.6). Il suffit donc de démontrer que les morphismes de
spécialisation sont injectifs (\hyperrefext[IX][IX.prop2.13]{2.13} (ii)). Pour cela, on
est réduit 
au cas $S$ strictement local, et par le théorème de changement de base
\hyperrefext[XII][XII.thm5.1]{5.1} (ii) on peut même remplacer $S$ par le spectre d'un anneau de
valuation discrète et strictement local. De plus, on peut supposer $X$
connexe. Puisque\pageoriginale $H^{1}(X,G)$ classifie les revêtements
principaux galoisiens de $X$ de groupe $G$, il suffit de vérifier que
si $X'\rightarrow X$ est un revêtement étale connexe alors la fibre
géométrique est également connexe. Cela résulte facilement du fait que
$f$ est lisse et propre, à fibres géométriques connexes.

Traitons maintenant le cas $F=(\ZZ/\ell)_X$. Pour les valeurs 0
et 1 de $q$ le résultat est déjà connu ((\hyperref[XIV.lem1.5]{1.5}) et
ci-dessus).  
\sisi{D'après  \hyperrefext[X][X.cor4.3]{4.3} et}{D'après} \hyperrefext[XII][XII.corbis5.3]{5.3 bis}, 
il reste seulement le cas
$q=2$, et de plus, si 
$S$ est de caractéristique $\ell$, on a
$R^{2}f_{\ast}(\ZZ/\ell)=0$. En remplaçant $S$ par un ouvert
non-vide convenable, on est donc réduit au cas où $\ell$ est
inversible sur $S$. Alors, d'après \hyperrefext[XII][XII.cor5.2]{5.2} \sisi{(ii)}{(iii)} 
et \sisi{X 4.7}{\hyperrefext[X][X.cor5.2]{5.2}}, la fibre de
$R^{2}f_{\ast}(\ZZ/\ell)$ en un point géométrique $\bar{s}$
au-dessus de $s \in S$ est isomorphe à $(\ZZ/\ell)^{c(s)}$, où
$c(s)$ est le nombre des composantes irréductibles de dimension 1 de
la fibre géométrique $X_{\bar{s}}$. Or $c$ est une fonction
constructible sur $S$ (\EGA IV 9.7.9), d'où le résultat (\hyperrefext[IX][IX.prop2.13]{2.13}
(iii)).

\begin{enonce*}[remark]{Remarque 1.11}\etiquette[1.11]{XIV.rem1.11}
%%%
%%% Since there is already a subsection 1.11, this remark should be numbered 1.12.
%%% I hesitate to change this, since it potentially affects references in other parts of SGA4.
%%% At least in XIV, there is no reference to Remark 1.11.
%%%
La démonstration de \hyperref[XIV.thm1.1]{1.1} se simplifie beaucoup en utilisant le
formalisme de la cohomologie à supports propres, qui sera développé
dans \ref{XVII} comme conséquence directe du théorème de changement de base
pour un morphisme propre \hyperrefext[XII][XII.thm5.1]{5.1}. D'autre part, l'énoncé
\hyperref[XIV.thm1.1]{1.1} se 
généralise au cas où on se donne sur $\sisi{Y}{S}$ un faisceau quelconque
d'anneaux de torsion $𝒜$, et qu'on prend sur $X$ un complexe
$K^{\bigdot}$ de $f^{\ast}(𝒜)$-Modules satisfaisant une condition
de « constructibilité » qui, dans le cas particulier \hyperref[XIV.thm1.1]{1.1},
s'exprimerait simplement en disant que des faisceaux de cohomologie
$\Hh^{i}(K^{\bigdot})$ de $K^{\bigdot}$ sont des $A$-modules
constructibles, nuls pour $i$ assez grand.
\end{enonce*}

\section{Une variante de la dimension}\etiquette[2]{XIV.sec2}\pageoriginale

\subsection{}\etiquette[2.1]{XIV.subsec2.1}
Soit $X$ un schéma. On  va noter
\begin{equation*}
d(x)=\dim\bar{\{x\}}\tag{2.1.1}
\end{equation*}\etiquette[2.1.1]{XIV.eq2.1.1}
pour $x\in X$, où $\bar{\{x\}}$ est l'adhérence de $\{x\}$. Si $F$ est
un faisceau abélien sur $X$, on pose
\begin{equation*}
d(F)=\sup \{d(x)\mid x\in X \quaet F_{\bar{x}}\neq 0\}{\quoi}.\tag{2.1.2}
\end{equation*}\etiquette[2.1.2]{XIV.eq2.1.2}
On aura besoin d'une notion rectifiée dans le cas local:

\begin{enonce*}{Définition 2.2}\etiquette[2.2]{XIV.def2.2}
Soient $Y$ le spectre d'un anneau local noethérien universellement
caténaire (\EGA IV 5.6.2) et $f:X\rightarrow Y$ un morphisme de
type fini. Soit $x\in X$ et $y=f(x)$. On pose
$$
\delta(x)=\dim\bar{\{y\}}+\degtr k(x)/k(y){\quoi}.
$$
Si $F$ est un faisceau abélien sur $X$, on pose
$$
\delta(F)=\sup \{\delta(x)\mid x\in X\quaet F_{\bar{x}}\neq 0\}{\quoi}.
$$
\end{enonce*}

\begin{enonce*}{Proposition 2.3}\etiquette[2.3]{XIV.prop2.3}
\begin{enumerate}
%%%
%%% The typesetting of this enumerate environment is terrible. 
%%% The spacing between the dash and (i) is too big. 
%%% Also, the items (ii), ... should better be without indent. 
%%% If you use the package enum.sty, then you can build the necessary list environment.
%%%
\renewcommand{\theenumi}{\roman{enumi}}
\renewcommand{\labelenumi}{(\theenumi)}
\item Soit $i:X'\rightarrow X$ une immersion et $x'\in X'$. Alors
$\delta(i(x'))=\delta(x')$.
%%% new line break within formula in text: spend extra line to induce new line breaks?

\item Soit $X_{0}$ la fibre fermée de $X/Y$. On a
$$
\delta(x)=
\begin{cases}
d(x) & \mbox{si } \bar{\{x\}}\cap X_{0}\neq ∅\\
d(x)+1 & \mbox{si } \bar{\{x\}}\cap X_{0}=∅{\quoi};
\end{cases}
$$

\item  $\delta(x)=d(x)$ si $f$ est propre.
\end{enumerate}
\end{enonce*}

\noindent
\emph{Démonstration}.\pageoriginale L'assertion (i) est triviale à partir de la
définition. Puisque $\bar{\{x\}}\cap X_{0}$ n'est pas vide dans le
cas $f$ propre, (iii) est conséquence de (ii). Pour (ii), nous pouvons
remplacer $X$ par $\bar{\{x\}}$ et $Y$ par $\bar{\{y\}}$ (avec la
structure induite réduite), ce qui est permis d'après (i). Soient
$x'\in X'$ et $y'=f(x')$. On a la formule (\EGA IV 5.6.1.1)
$$
\dim 𝒪_{X,x'}=\dim
𝒪_{Y,y'}+\degtr k(x)/k(y) - \degtr k(x')/k(y'){\quoi}.
$$

Supposons que $X_{0}\neq ∅$. Alors il est clair que 
$\displaystyle{\mathop{\sup}_{x'}}$
$\{\dim 𝒪_{X,x'}\}=\dim X$ est obtenu en prenant pour $x'$ un point
fermé de $X_{0}$, d'où
$$
\dim X=\dim 𝒪_{X,x'}=\dim Y + \degtr k(x)/k(y)=\delta(x){\quoi}.
$$
Si $X_{0}=∅$, le supremum est évidemment obtenu en prenant pour
$x'$ un point fermé dans sa fibre $f^{-1}(y')$, tel que
$\dim\bar{\{y'\}}=1$. Il existe de tels points, parce que les $y'\in
Y$ avec $\dim\bar{\{y'\}}=1$ forment un ensemble très dense (\EGA IV
10.1.3, \sisi{10.5.9, 1.8.4}{10.3.1, 10.4.1, 10.5.9}). On a donc:
\begin{align*}
\dim X=\dim 𝒪_{X,x'} &= \dim 𝒪_{Y,y'}+\degtr k(x)/k(y)\\
&= \dim Y-1 +\degtr k(x)/k(y)=\delta(x) \sisi{=}{-} 1{\quoi},
\end{align*}
d'où le résultat.

\begin{enonce*}{Proposition 2.4}\etiquette[2.4]{XIV.prop2.4}
Soit $f:X\rightarrow Y$ un morphisme de schémas de type fini sur le
spectre d'un corps $k$. Soient $x\in X$, $y=f(x)$, et $y_{0}$ une
spécialisation de $y$. Soient $f':X'\rightarrow Y'$ le localisé strict
de $f$ en un point géométrique $\bar{y}_{0}$ au-dessus de
$y_{0}$, et $x'\in X'$ un point au-dessus de $x$. Alors on a
$$
\delta(x')=d(x)-d(y_{0}){\quoi}.
$$
\end{enonce*}

\noindent
\emph{Démonstration}.\pageoriginale On a
\begin{align*}
d(x) &= d(y)+\degtr k(x)/k(y)\\
\delta(x') &= d(y')+\degtr k(x')/k(y'){\quoi},
\end{align*}
où $y'=f'(x')$. De plus
$$
\degtr k(x)/k(y)=\degtr k(x')/k(y'){\quoi},
$$
parce que les extensions $k(x')/k(x)$ et $k(y')/k(y)$ sont
algébriques. Il reste donc à démontrer que
$$
d(y')=d(y)-d(y_{0}){\quoi}.
$$

On peut supposer $y$ maximal dans $Y$, donc $y'$ maximal dans $Y'$, et
alors
\begin{align*}
d(y') &= \dim \Oo_{Y',\bar{y}_{0}}\\
      &= \dim \Oo_{Y,y_{0}}\quad \text{\rm (\EGA IV 6.1.3)}\\
      &= \sisi{\dim Y}{\dim_{y_0} Y} -d(y_{0})\quad \text{\rm (\EGA IV 5.2.3.1)}\\
      &= d(y)-d(y_{0}){\quoi},
\end{align*}
d'où le résultat.

\section{Dimension cohomologique des schémas algébriques affines}\etiquette[3]{XIV.sec3}

\begin{enonce*}{Théorème 3.1~(d)}\etiquette[3.1]{XIV.thm3.1}\nde{Le théorème \ref{XIV.thm1.1} est aujourd'hui connu sous le nom de \textit{théorème de Lefschetz affine}.}
Soit $f:X\rightarrow Y$ un morphisme    affine de schémas de type fini
sur un corps $k$. Soit $F$ un faisceau de torsion sur $X$ tel que
$d(F)\leq d$ \emph{(\cf \hyperref[XIV.subsec2.1]{2.1})}. 
Alors $d(R^{q}f_{\ast}F)\leq d-q$.
\end{enonce*}

En particulier, prenant $Y=\Spec(k)$:

\begin{enonce*}{Corollaire 3.2}\etiquette[3.2]{XIV.cor3.2}
Soit\pageoriginale $X$ un schéma affine de type fini sur un corps $k$
séparablement clos. 
Alors\nde{Considérons l'assertion analogue en géométrie analytique complexe:
Soit $X$ une variété affine sur $\mathbf{C}$ et soit $F$ un faisceau constructible sur $X$. Alors, on a $H^m(X^{\rm an},F) = 0$ pour tout $m>\dim X$, où $X^{\rm an}$ est l'espace analytique associée à $X$. Bien que cette assertion découle de \hyperref[XIV.cor3.2]{3.2} par le théorème de comparaison \hyperrefext[XVI][XVI.thm4.1]{4.1}, on a aussi deux preuves entièrement de nature analytique données par Pierre Deligne et Hélène Esnault, voir Esnault, H., Variation on Artin's vanishing theorem, \textit{Advances in Math.} \textbf{198} (2005), 435--438.}
$$
\cd X\leq \dim X{\quoi}.
$$
\end{enonce*}

\begin{enonce*}{Corollaire 3.3}\etiquette[3.3]{XIV.cor3.3}
Soit $X\subset \PP^N_k$ un schéma projectif de dimension $n$
sur un corps séparablement clos $k$, et soit $Y=H\cap X$ une « section hyperplane »  de $X$. Alors pour tout faisceau de torsion
$F$ sur $X$, l'homomorphisme canonique
$$
H^{q}_Y(X,F)\longrightarrow H^{q}(X,F)
$$
est bijectif pour $q>n+1$, et surjectif pour $q=n+1$.
\end{enonce*}

Le corollaire est trivial à partir de \sisi{\hyperref[XIV.def2.2]{2.2}}{\hyperref[XIV.cor3.2]{3.2}} et de la suite exacte pour
un sous-ensemble fermé (\sisi{IV 3.7.1}{\hyperrefext[V][V.6.5.4]{6.5.4}}), 
%%% This last reference should read \ref{V.eq6.5.4}, but the TeX-file of chapter V contains the
%%% label without "eq".
compte tenu que $X-Y=\sisi{\mathcal{U}}{U}$ est
affine de dimension $\leq n$. Remarquons que \hyperref[XIV.cor3.3]{3.3} est une
généralisation d'un des théorèmes de Lefschetz sur les sections
hyperplanes: en effet, si $X$ est lisse sur $k$ et $F$ est localement
constant et premier à la caractéristique, on peut calculer la
cohomologie $H^{q}_Y(X,F)$ explicitement en utilisant le théorème de
pureté cohomologique relatif (\hyperrefext[XVI][XVI.sec3]{3}) et la suite spectrale (\sisi{\quad}{\hyperrefext[V][V.6.4.3]{6.4.3}}), 
%%% This last reference should read \ref{V.eq6.4.3}, but the TeX-file of chapter V contains the
%%% label without "eq".
et
on trouve que
$$
H^{q}_Y(X,\ZZ/n)\simeq H^{q-2}(Y,\mu_n^{-1}) \mbox{\rm\quad pour
tout $q$;}
$$
l'homomorphisme $H^{q-2}$ $(Y,\mu_n^{-1})\rightarrow
H^{q}(X,\ZZ/n)$ déduit de \hyperref[XIV.cor3.3]{3.3} n'étant autre que 
l'~« homomorphisme de Gysin », qui sera étudié ultérieurement
\footnote[1]{\Cf \SGA 2 XIV pour une étude systématique des théorèmes
du type de Lefschetz.}.

Le théorème \hyperref[XIV.thm3.1]{3.1}~(d) équivaut\nde{Plus précisément,
l'argument du texte donne le résultat suivant. L'assertion \hyperref[XIV.cor3.4]{3.4}~(d') pour $d' \leq d$ implique \hyperref[XIV.thm3.1]{3.1}~(d), et l'assertion \hyperref[XIV.thm3.1]{3.1}~(d) implique \hyperref[XIV.cor3.4]{3.4}~(d). Cela suffit pour conclure ci-après.}
 au suivant:

\begin{enonce*}{Corollaire 3.4~(d)}\etiquette[3.4]{XIV.cor3.4}
Soit $Y'$ un schéma strictement local qui est un localisé strict d'un
schéma algébrique sur un corps. Soit $f':X'\rightarrow Y'$ un
morphisme affine\pageoriginale de type fini, et $F'$ un faisceau de
torsion sur $X'$ tel que $\delta(F')\leq d$ \emph{(\cf (\hyperref[XIV.def2.2]{2.2}))}. Alors
$H^{q}(X',F')=0$ pour $q>d$.
\end{enonce*}

En effet, supposons que \hyperref[XIV.cor3.4]{3.4}~(d) soit vrai. Soit $f:X\rightarrow Y$ un
morphisme affine de schémas algébriques sur $k$ et soit $F$ un
faisceau de torsion sur $X$ tel que $d(F)\leq d$. Soit $y$ un point de
$Y$ et $\xi$ un point géométrique au-dessus de $y$. Pour démontrer 
\sisi{\ref{XIV.thm3.1}}{\hyperref[XIV.thm3.1]{3.1}} (d), 
il faut démontrer que la fibre $(R^{q}f_{\ast}F)_{\xi}$ est nulle
si $d(y)>d-q$.

Soit $f':X'\rightarrow Y'$ le localisé strict de $f$ au point $\xi$,
c'est-à-dire, soit $Y'$ le localisé strict de $Y$ en $\xi$ et soit
$X'=\fprod{X}{Y'}{Y}$.

Soit $F'$ le faisceau induit de $F$ sur $X$, de sorte que
$(R^{q}f_{\ast}F)_{\xi}=H^{q}(X',F')$ (\hyperrefext[VIII][VIII.thm5.2]{5.2}). 
\sisi{On a
$$
\delta(F')\leq d(F)-d(y)\leq d-d(y)\quad (\hyperref[XIV.prop2.4]{2.4}){\quoi},
$$}
{On a d'après (\hyperref[XIV.prop2.4]{2.4})
$$
\delta(F')\leq d(F)-d(y)\leq d-d(y){\quoi},
$$}
%%% I find it irritating that the reference to the proposition from which the above follows,
%%% is referred to in a style that might be confused with an equation number 
%%% (although the correct equation number is 3.4.1)
et en appliquant \hyperref[XIV.cor3.4]{3.4}~(d), on trouve $H^{q}(X',F')=0$ si $q>d-d(y)$,
d'où le résultat.

Inversement, supposons que \sisi{\hyperref[XIV.subsec2.1]{2.1}}{\hyperref[XIV.thm3.1]{3.1}}~(d) 
soit vrai, et soient $F'$,
$f':X'\rightarrow Y'$ comme dans l'énoncé de \hyperref[XIV.cor3.4]{3.4}~(d). Puisque la
cohomologie commute aux limites inductives (\hyperrefext[VII][VII.prop3.3]{3.3}) et $F'$ est
limite inductive de ses sous-faisceaux constructibles (\hyperrefext[IX][IX.prop2.9]{2.9}), on
peut supposer $F'$ constructible. Or $Y'$ est localisé strict d'un
schéma algébrique $Y$ en un point géométrique $\xi$ au-dessus d'un
point \emph{fermé} $y_{0}$ (\hyperrefext[X][X.lem3.3]{3.3}). Écrivons
$Y'=\varprojlim_{\alpha} Y_{\alpha}$, où $(Y_{\alpha})$ est un système
projectif filtrant de schémas affines étales sur $Y$ 
(\hyperrefext[VIII][VIII.subsec4.5]{4.5}). D'après \hyperrefext[IX][IX.cor2.7.4]{2.7.4}, $F'$ est induit d'un faisceau de torsion
constructible \sisi{}{$F_\alpha$} sur 
$\sisi{X=\fprod{X}{Y_K}{Y}}{X_\alpha = \fprod{X}{Y_\alpha}{Y}}$ pour $\alpha$ suffisamment
grand, et on voit immédiatement, utilisant 2.4 et\pageoriginale \EGA IV
9.5.5, qu'on peut de plus supposer que $d(F_{\alpha})=\delta(F')\leq
d$. En remplaçant $Y$ par $Y_{\alpha}$, on est réduit au cas où
$F'$ est induit par un faisceau $F$ sur $X$, avec $d(F)\leq d$. On a
alors $R^{q}f_{\ast}F=0$ si $q>\sisi{D}{d}$, 
d'après \sisi{\hyperref[XIV.subsec2.1]{2.1}}{\hyperref[XIV.thm3.1]{3.1}}~(d), donc en
particulier
$$
(R^{q}f_{\ast}F)_{\xi}=H^{q}(X',F')=0\quad{\mbox{\rm si}}\quad
q>d{\quoi},
$$
d'où \hyperref[XIV.cor3.4]{3.4}~(d).
 
Comme cas particulier de \hyperref[XIV.cor3.4]{3.4}~(d), on a

\begin{enonce*}{Corollaire 3.5}\etiquette[3.5]{XIV.cor3.5}
Soit $Y'$ un schéma strictement local de dimension $\leq d$, localisé
strict d'un schéma algébrique sur un corps $k$. Soit $U\subset Y'$ un
ouvert affine. Alors $\cd U\leq d$.
\end{enonce*}

\begin{enonce*}[remark]{Remarque 3.6}\etiquette[3.6]{XIV.rem3.6}
Il semble très plausible\nde{Ofer Gabber a réussit à démontrer l'analogue suivant de \ref{XIV.thm3.1}: Soit $f: X \to Y$ un morphisme affine de type fini avec $Y$ quasi-excellent (\EGA IV 7.8) possédant une fonction de dimension $\delta_Y$. Soit $\delta_X$ la fonction de dimension rectifiée associée sur $X$. Pour $n$ inversible sur $Y$ et tout faisceau $F$ de $\mathbf{Z}/n$ sur $X$ on a $\delta_Y(R^q f_\ast F) \leq \delta_X(F) - q$ pour tout $q \geq 0$. Ici, on définit comme ci-dessus $\delta(G) = \sup \{\delta(x) \mid x\in X \quaet G_{\bar{x}}\neq 0\}$ pour un faisceau $G$.
} 
que \hyperref[XIV.cor3.4]{3.4}, donc \hyperref[XIV.cor3.5]{3.5}, 
reste valable pour
\emph{tout} schéma local noethérien $Y'$ (pas nécessairement un
localisé strict d'un schéma algébrique), du moins si $Y'$ est
\emph{excellent} (\EGA IV 7.8). C'est ce qui sera prouvé dans (\sisi{XX 6}{\hyperrefext[XIX][XIX.sec6]{6}})
lorsqu'on suppose de plus $Y$ de caractéristique nulle.
\end{enonce*}

\section{Démonstration du théorème 3.1}\etiquette[4]{XIV.sec4}

La démonstration se fait par récurrence sur $d$.

\begin{enonce*}{Lemme 4.1}\etiquette[4.1]{XIV.lem4.1}
\emph{\hyperref[XIV.thm3.1]{3.1}~(0)} et \emph{\hyperref[XIV.thm3.1]{3.1}~(1)} sont vrais.
\end{enonce*}

\noindent
\emph{Démonstration}. Triviale pour $d=0$, auquel cas $\supp F$ est
fini sur $k$, et on applique (\hyperrefext[VIII][VIII.prop5.5]{5.5}). Soient $F$, $f:X\rightarrow
Y$ comme dans l'énoncé de \hyperref[XIV.thm3.1]{3.1} et supposons que $d(F)\leq 1$. En
remplaçant $F$ par un sous-faisceau constructible (\hyperrefext[IX][IX.prop2.9]{2.9} et \hyperrefext[VII][VII.prop3.3]{3.3}), 
on se réduit au cas où $F$ est constructible. Alors le
support\pageoriginale de $F$ est contenu dans un sous-schéma fermé de
$X$ de dimension $\leq 1$ (\hyperrefext[IX][IX.def2.3]{2.3}\sisi{ (ii)}{}), et on peut remplacer $X$ par
ce sous-schéma, c'est-à-dire, on peut supposer qu'on a $\dim X\leq
1$. D'après \hyperrefext[IX][IX.cor5.7]{5.7}, la dimension cohomologique d'un schéma algébrique
affine de dimension $1$ sur un corps séparablement clos est $\leq 1$. Il
suffit d'ailleurs de traiter le cas où $k$ est séparablement clos: en
effet, soit $\bar{f}:\bar{X}\rightarrow \bar{Y}$ le morphisme déduit
de $f$ par le changement de base $\Spec \bar{k}\rightarrow \Spec k$
($\bar{k}$ la clôture séparable de $k$). Pour un point géométrique
$\xi$ de $Y$, le localisé strict $f':X'\rightarrow Y'$ de $f$ en $\xi$
s'identifie à un localisé strict de $\bar{f}$, donc la fibre
$(R^{q}f_{\ast}F)_{\xi}$ à une fibre de
$R^{q}\bar{f}_{\ast}\bar{F}$. Puisque $f$ est affine, on trouve donc
que $R^{q}f_{\ast}F=0$ si $q>1$. De plus, il est évident que
$d(f_{\ast}F)\leq 1$. Pour le cas $q=1$, notons que
l'hypothèse que $X$ soit de dimension~$\leq 1$ implique qu'il existe
une partie fermée finie $Z$ de $Y$ telle que $f$ soit fini au-dessus
de $Y-Z$. Il en résulte que $R^{1}f_{\ast}F$ est nul dans $Y-Z$ (\hyperrefext[VIII][VIII.cor5.6]{5.6}), 
d'où $d(R^{1}f_{\ast}F)\leq 1$, ce qui achève la démonstration
du lemme.

Soit $d\geq 2$ et supposons maintenant que \hyperref[XIV.thm3.1]{3.1}~(d') (ou, ce
qui revient au même, \hyperref[XIV.cor3.4]{3.4}~(d')) est vrai pour $d'<d$.

\begin{enonce*}{Lemme 4.2}\etiquette[4.2]{XIV.lem4.2}
Pour démontrer \emph{\hyperref[XIV.cor3.4]{3.4}~(d)} dans tous les cas, il suffit de traiter
le cas où $X'=\EE^{1}_{Y'}$ est l'espace affine dimension $1$
sur $Y'$, et où $f$ est le morphisme structural.
\end{enonce*}

\noindent
\emph{Démonstration}. Évidemment, on peut supposer
$X'=\EE^{N}_{Y'}$ pour $n$ convenable, parce qu'on peut plonger
$X'$ dans un $\EE^{N}_{Y'}$. Supposons $N>1$ et le résultat
connu pour tout $Y'$ et pour $\EE^{r}_{Y'}$ avec $r<N$. Soit
$F'$ un faisceau de torsion sur $\EE^{N}_{Y'}$,
avec\pageoriginale $d(F')\leq d$. Considérons le diagramme
$$
\xymatrix@=1.5cm{
\EE^{N}_{Y'} \ar[rd]_{f'}       \ar[rr]^{g'}&      &
\EE^{N-1}_{Y'} \ar[ld]^{h'}\\
                       & Y' & &,
}
$$
$g'$ et $h'$ les projections canoniques, de sorte que
$\EE^{N}_{Y'}$ est isomorphe à l'espace affine de dimension $1$
au-dessus de $\EE^{N-1}_{Y'}$. Puisque évidemment chaque anneau
localisé strict de $\EE^{N-1}_{Y'}$ est aussi le localisé
strict d'un schéma algébrique, on peut appliquer l'hypothèse de
récurrence aux fibres de $R^{q}g'_{\ast}F'$, et on trouve que
\begin{equation*}
\delta(R^{q}g'_{\ast}F')\leq d-q{\quoi}.\tag{$\ast$}
\end{equation*}
De plus, \hyperref[XIV.cor3.4]{3.4}~(d') est vrai pour $h'$, $d'\leq d$, par
l'hypothèse de récurrence sur $N$.

Considérons la suite spectrale de Leray
$$
E^{p,q}_2=H^{p}(\EE^{N-1}_{Y'}, R^{q}
g'_{\ast}(F'))\Longrightarrow H^{n}(\EE^{N}_{Y'}, F'){\quoi}.
$$
On a $E^{p,q}_2=0$ si $p>d-q$, comme on trouve en appliquant $(\ast)$
et \hyperref[XIV.cor3.4]{3.4}~(d') au morphisme $h'$. Il s'ensuit que l'aboutissement
est nul si $n>d$, d'où le lemme.

Considérons l'énoncé supplémentaire suivant, qui est un cas spécial de \hyperref[XIV.cor3.5]{3.5}:

\begin{enonce*}{Énoncé 4.3~(d)}\etiquette[4.3]{XIV.enon4.3}
Soit $X'$ strictement local de dimension $\leq d$, localisé strict
d'un schéma algébrique sur un corps. Soit $f\in \Gamma$
$(X',\Oo_{X'})$ et soit $U\subset X'$ l'ouvert $U=X'-V(f)$. 
%%% new line break within formula in text: spend extra line to induce new line breaks?
Alors $\cd U\leq d$.
\end{enonce*}

\begin{enonce*}{Lemme 4.4}\etiquette[4.4]{XIV.lem4.4}
\emph{\hyperref[XIV.enon4.3]{4.3}~(d)} implique \emph{\hyperref[XIV.cor3.4]{3.4}~(d)}.
\end{enonce*}

\noindent
\emph{Démonstration}.\pageoriginale Considérons le diagramme
$$
\xymatrix@=1.5cm{
\EE^{1}_{Y'} \ar[rr]^{i} \ar[rd]_{f'}&            &
\PP^{1}_{Y'} \ar[ld]^{g'}\\
                    &       Y'   & 
}
$$
où $Y'$ est comme dans l'énoncé de \sisi{2.6}{\hyperref[XIV.cor3.4]{3.4}} et où $\PP^{1}_{Y'}$ est
l'espace projectif, déduit de $\EE^{1}$ « en ajoutant la
section $Y'_{\infty}$ à l'infini ». Soit $F'$ un faisceau de torsion
sur $\EE^{1}$ avec $d(\supp F')\leq d$. On a la suite spectrale
$$
H^{p}(\PP^{1}, R^{q}i_{\ast}F')\Longrightarrow
H^{n}(\EE^{1},F'){\quoi}.
$$

Or les $R^{q}i_{\ast}F$, pour $q>0$, sont concentrés sur
$Y'_{\infty}$, qui est $Y'$-isomorphe à $Y'$. Puisque $Y'$ est
strictement local, il s'ensuit que $H^{p}$ $(\PP^{1},
R^{q}i_{\ast}F')=0$ si $p$, $q>0$. De plus, d'après le \emph{théorème
de changement de base pour un morphisme propre} (\hyperrefext[XII][XII.thm5.1]{5.1}), la
cohomologie d'un faisceau de torsion sur $\PP^{1}$ peut-être
calculée sur la fibre fermée, qui est un schéma de dimension
cohomologique $\leq 2$ (\hyperrefext[X][X.cor4.3]{4.3}). On a donc $H^{p}$ $(\PP^{1},
i_{\ast}F')=0$, si $p>2$. Comme on veut démontrer $H^{n}
(\EE^{1},F')=0$ pour $n>d$, et comme $d\geq 2$, on est réduit
à démontrer que $H^{0} (\PP^{1}, R^{q}i_{\ast}F')=0$ pour
$q>d$. Ce groupe, isomorphe aussi à $H^{0} (Y'_{\infty},
R^{\sisi{1}{q}}i_{\ast}(F'))$, est en vertu de (\hyperrefext[VIII][VIII.prop4.6]{4.6}) isomorphe à la fibre
de $R^{q}i_{\ast}F'$ au point $Q$, $Q$ étant le point à l'infini de
$\PP^{1}$ dans la fibre fermée.

Il suffit (\hyperrefext[VII][VII.prop3.3]{3.3} et \hyperrefext[IX][IX.prop2.9]{2.9}) de traiter le cas où $F'$ est
constructible. Alors, puisque $\delta(\supp F')\leq d$, il existe un
sous-schéma fermé $X'$ de $\PP^{1}$ avec $\dim X'\leq d$, tel
que $X'$ contienne le support de $F'$. Soit $\widetilde{X}'$ le
localisé de $X'$ au point géométrique $Q$, soit
$\widetilde{U}=\widetilde{X}'-
\fprod{\widetilde{X}}{Y'_{\infty}}{\PP^{1}}$, et soit $F$ le
faisceau\pageoriginale induit par $F'$ sur $U$ par le morphisme
$\widetilde{U}\rightarrow \PP^{1}$. Alors la fibre de
$R^{q}i_{\ast}F$ au point $Q$ n'est autre que $H^{\sisi{p}{q}}(\widetilde{U},
\widetilde{F})$ (\hyperrefext[VIII][VIII.thm5.2]{5.2}), et on applique \hyperref[XIV.enon4.3]{4.3}~(d) au schéma
strictement local $X'$ et à l'ouvert $\widetilde{U}$.

Le lemme ci-dessous, joint à \hyperref[XIV.lem4.4]{4.4}, achèvera la démonstration.

\begin{enonce*}{Lemme 4.5}\etiquette[4.5]{XIV.lem4.5}
\emph{\hyperref[XIV.cor3.4]{3.4}~ (d')} pour $d'<d$ implique
\emph{\hyperref[XIV.enon4.3]{4.3}~(d')} 
pour $d'\leq d$.
\end{enonce*}

\noindent
\emph{Démonstration}. Supposons \hyperref[XIV.cor3.4]{3.4}~(d') pour $d'<d$. Par
récurrence, nous pouvons supposer que \hyperref[XIV.enon4.3]{4.3}~(d') est vrai pour
$d'<d$. Soient les données $X'$, $f$, $U$ comme dans l'énoncé \hyperref[XIV.enon4.3]{4.3}~(d). 
Alors $X$ est localisé strict d'un schéma algébrique $X_{0}$
sur un corps $k$ en un point géométrique \emph{fermé}. On peut prendre
$X_{0}$ affine et de dimension $\leq d$, quitte à le remplacer par
un voisinage de $x_{0}$. Alors $X$ est limite projective de
schémas affines et étales $X_i$ au-dessus de $X_{0}$ (\hyperrefext[VIII][VIII.subsec4.5]{4.5}).

Soit $F$ un faisceau de torsion sur $U$. Il faut démontrer que
$H^{q}(U,F)=0$ si $q>d$, et il suffit encore de traiter le cas où $F$
est constructible. Alors les données $f$, $F$, $U$ proviennent de
données analogues sur l'un des $X_i$ (\sisi{\hyperrefext[IX][IX.prop2.7]{2.7}}{\hyperrefext[IX][IX.cor2.7.4]{2.7.4}}). 
On aura un faisceau
$F_i$ sur $U_i=X_i-V(f_i)$ qui « induit »  $F$ sur $U$. Puisque
$U=\varprojlim U_i$, on a (\hyperrefext[VII][VII.cor5.8]{5.8})
$$
H^{q}(U,F)=\varinjlim H^{q}(U_i, F_i){\quoi}.
$$
Il suffit donc de démontrer le fait suivant: Quel que soit $i$, il
existe un diagramme commutatif
$$
\xymatrix@=1.5cm{
 U \ar[rr] \ar[rd]&        & U''\ar[ld]\\
   &  U_i   & 
}
$$
où $U\rightarrow U_i$ est la flèche donnée, tel que $\cd(U'')\sisi{<}{\leq}d$. En
effet, cela impliquera que\pageoriginale chaque élément de
$H^{q}(U_i,F_i)$, pour $q>d$, a une image nulle dans $H^{q}(U,F)$.

Supposons $i=0$, et considérons le morphisme $g:X_{0}\rightarrow
\EE^{1}_k=Y_{0}$ donnée par $f_{0}\in\Gamma$
$(X_{0},\Oo_{X_{0}})$. Soit $Y''$ le localisé strict de
$\EE^{1}_k$ en un point géométrique au-dessus de l'origine,
donc $Y''$ est le spectre d'un anneau de valuation discrète, et soit
$X''=\fprod{X_{0}}{Y''}{Y_{0}}$.
%%% new line break within formula in text: spend extra line to induce new line breaks?


Le morphisme $X\rightarrow X_{0}$ se factorise par $X''$ puisque
$X$ est strictement local. De plus, l'ouvert
$U''=\fprod{X''}{U_{0}}{X_{0}}$ n'est autre que la fibre
générique de $X''/Y''$. C'est donc un schéma algébrique affine
au-dessus du point générique de $Y'$, et on a évidemment $\dim U''\leq
d-1$.

Vérifions que $\cd U''\leq d$ (ce qui achèvera la démonstration). Soit
$K$ le corps résiduel de $Y''$ au point générique. C'est un corps de
dimension cohomologique 1, c'est-à-dire, on a $\cd(G)=1$, $G$ le
groupe de Galois de la clôture séparable $\bar{K}$ de $K$ 
(\hyperrefext[X][X.thm2.2]{2.2}). Soit $\bar{U}''=\fprod{U''}{(\Spec \bar{K})}{\Spec K}$. On a
$\cd\bar{U}''\leq d-1$ par hypothèse de récurrence sur $d$ (c'est
\sisi{\hyperref[XIV.def2.2]{2.2} pour un schéma de dimension $\leq d$}
{\hyperref[XIV.cor3.2]{3.2} pour un schéma de dimension $< d$}). Le fait que $\cd U''\leq d$
suit alors de la suite spectrale de Hochschild-Serre (\hyperrefext[VIII][VIII.prop8.4]{8.4})
$$
H^{q}(G,H^{q}(\bar{U}\sisi{'}{''},\bar{F}))\Longrightarrow H^{n}(U,F){\quoi},
$$
d'où le lemme \hyperref[XIV.lem4.5]{4.5}.

\ifx\danslelivre\undefined
\end{document}
\fi
